\documentclass[11pt,a4paper]{article}

\usepackage[utf8]{inputenc}
\usepackage[T1]{fontenc}
\usepackage[brazil]{babel}
\usepackage{geometry}
\usepackage{xcolor}
\usepackage{enumitem}
\usepackage{listings}
\usepackage{hyperref}

\geometry{margin=1.65cm}
\setlength{\parindent}{0pt}
\setlength{\parskip}{4pt}

\definecolor{rankblue}{RGB}{28,72,132}
\definecolor{softgray}{RGB}{245,246,248}

\lstset{
  basicstyle=\ttfamily\footnotesize,
  breaklines=true,
  frame=single,
  columns=fullflexible,
  keepspaces=true,
  showstringspaces=false,
  backgroundcolor=\color{softgray}
}

\title{\vspace{-1.2cm}20 Perguntas Mais Prováveis da Defesa \\ Simulador TR1}
\author{}
\date{}

\newcommand{\rank}[1]{\textcolor{rankblue}{\textbf{Rank #1}}}
\newcommand{\pergunta}[2]{\item[\rank{#1}.] #2}

\begin{document}

\maketitle
\vspace{-1.0cm}

\section*{Perguntas ranqueadas por probabilidade}

\begin{enumerate}[leftmargin=2.3cm, itemsep=2pt]
  \pergunta{1}{Como vocês modulam com ruído?}
  \pergunta{2}{Como foi feita a modulação no projeto?}
  \pergunta{3}{O ruído entra antes ou depois da modulação? Em qual arquivo isso acontece?}
  \pergunta{4}{Qual é o jeito correto de colocar o bit de paridade? Ele fica no final de cada byte?}
  \pergunta{5}{Por que usar bit de paridade se o CRC-32 detecta melhor?}
  \pergunta{6}{Qual é a vantagem do Manchester em relação ao NRZ-Polar e ao Bipolar AMI?}
  \pergunta{7}{Quais são as vantagens e desvantagens de NRZ-Polar e 16-QAM?}
  \pergunta{8}{Explique o Hamming(8,4) implementado no trabalho.}
  \pergunta{9}{Por que o Hamming vem depois do EDC no transmissor e antes da verificação do EDC no receptor?}
  \pergunta{10}{Como vocês testaram o projeto e como provaram que o CRC-32 está correto?}
  \pergunta{11}{Como funciona o enquadramento com flags, byte stuffing e bit stuffing?}
  \pergunta{12}{Qual é a diferença entre paridade, checksum e CRC-32?}
  \pergunta{13}{Por que 16-QAM é mais sensível ao ruído do que QPSK?}
  \pergunta{14}{O que significa constelação I/Q e para que serve a função de correlação?}
  \pergunta{15}{Como ASK, FSK, QPSK e 16-QAM funcionam especificamente no código?}
  \pergunta{16}{O que acontece se o ruído aumentar muito?}
  \pergunta{17}{Por que o código completa bits com zero em QPSK e 16-QAM? Isso altera a mensagem?}
  \pergunta{18}{Como o receptor decide se recebeu 0 ou 1 nas modulações banda-base?}
  \pergunta{19}{Por que o projeto separa camada física, meio de comunicação e camada de enlace em arquivos diferentes?}
  \pergunta{20}{Se o professor pedir para abrir o código, quais funções você deve saber explicar primeiro?}
\end{enumerate}

\newpage

\section*{Gabarito e explicações}

\begin{enumerate}[leftmargin=1.2cm, itemsep=8pt]
  \item[\rank{1}.] \textbf{Como vocês modulam com ruído?}

  A resposta correta é: \textbf{não modulamos o ruído}. O projeto modula os bits, gera um sinal limpo e só depois o meio de comunicação adiciona ruído gaussiano ao sinal.

  Frase pronta:

  \begin{quote}
  Nós não modulamos com ruído. Primeiro a camada física gera o sinal limpo a partir dos bits. Depois o meio de comunicação soma ruído gaussiano em cada amostra. O receptor demodula o sinal já ruidoso.
  \end{quote}

  \item[\rank{2}.] \textbf{Como foi feita a modulação no projeto?}

  Existem dois grupos. Em banda-base, os bits viram níveis ou transições de tensão: NRZ-Polar, Manchester e Bipolar AMI. Em portadora, os bits escolhem propriedades de uma onda senoidal: ASK muda amplitude, FSK muda frequência, QPSK usa 2 bits por símbolo e 16-QAM usa 4 bits por símbolo.

  \begin{lstlisting}[language=Python]
amostras.append(i * math.cos(angulo) - q * math.sin(angulo))
  \end{lstlisting}

  \item[\rank{3}.] \textbf{O ruído entra antes ou depois da modulação? Em qual arquivo isso acontece?}

  O ruído entra \textbf{depois} da modulação. A ordem é:

  \begin{quote}
  bits -> modulação -> sinal limpo -> meio ruidoso -> sinal recebido -> demodulação
  \end{quote}

  Isso acontece no \texttt{meio\_comunicacao.py}, não dentro da camada física.

  \begin{lstlisting}[language=Python]
ruido = random.gauss(media, sigma)
saida.append(amostra + ruido)
  \end{lstlisting}

  \item[\rank{4}.] \textbf{Qual é o jeito correto de colocar o bit de paridade? Ele fica no final de cada byte?}

  Não. No trabalho, a paridade é calculada sobre o \textbf{quadro inteiro}, não no final de cada byte. O código soma todos os bits do bloco e calcula a paridade par. Como o simulador trabalha o EDC em bytes, a paridade ocupa 1 byte: sete zeros e o bit de paridade no final.

  \begin{lstlisting}[language=Python]
paridade = sum(bits) % 2
return bits + [0] * 7 + [paridade]
  \end{lstlisting}

  Frase pronta:

  \begin{quote}
  A paridade é por quadro. No nosso código ela ocupa um byte de EDC, mas só o último bit desse byte é a paridade.
  \end{quote}

  \item[\rank{5}.] \textbf{Por que usar bit de paridade se o CRC-32 detecta melhor?}

  Porque a paridade é simples e tem menor overhead. O CRC-32 é mais forte, principalmente para erros em rajada, mas adiciona 32 bits por quadro. A paridade, conceitualmente, precisa de 1 bit; no simulador, ocupa 1 byte por alinhamento do EDC.

  \begin{quote}
  Paridade é economia e simplicidade. CRC-32 é robustez. A escolha depende do custo aceito em bits extras e do nível de detecção desejado.
  \end{quote}

  \item[\rank{6}.] \textbf{Qual é a vantagem do Manchester em relação ao NRZ-Polar e ao Bipolar AMI?}

  Manchester sempre tem uma transição no meio do bit. Essa transição ajuda no sincronismo, porque o receptor tem uma mudança previsível dentro de cada intervalo de bit. No projeto, bit 0 é baixo para alto e bit 1 é alto para baixo.

  \begin{quote}
  A vantagem é sincronismo: o receptor não depende de longas sequências paradas, porque todo bit tem transição no meio.
  \end{quote}

  \item[\rank{7}.] \textbf{Quais são as vantagens e desvantagens de NRZ-Polar e 16-QAM?}

  NRZ-Polar é simples: bit 1 vira tensão positiva e bit 0 vira tensão negativa. É fácil de implementar e tende a ser mais robusto, mas carrega só 1 bit por símbolo e pode ter longas sequências sem transição.

  16-QAM é mais eficiente: carrega 4 bits por símbolo em uma grade I/Q. A desvantagem é que os pontos ficam mais próximos, então o ruído pode deslocar o símbolo para perto do ponto errado.

  \item[\rank{8}.] \textbf{Explique o Hamming(8,4) implementado no trabalho.}

  O Hamming(8,4) pega 4 bits de dados e gera 8 bits transmitidos. Ele usa três paridades de Hamming e uma paridade geral:

  \begin{quote}
  p1 p2 d1 p4 d2 d3 d4 p0
  \end{quote}

  A síndrome indica a posição de erro simples. A paridade geral permite diferenciar erro simples de erro duplo.

  \begin{lstlisting}[language=Python]
p1 = (d1 + d2 + d4) % 2
p2 = (d1 + d3 + d4) % 2
p4 = (d2 + d3 + d4) % 2
bloco = [p1, p2, d1, p4, d2, d3, d4]
p0 = sum(bloco) % 2
  \end{lstlisting}

  \item[\rank{9}.] \textbf{Por que o Hamming vem depois do EDC no transmissor e antes da verificação do EDC no receptor?}

  No transmissor, o EDC entra primeiro. Depois o Hamming protege tanto o payload quanto os bits de detecção. No receptor, o caminho é inverso: primeiro tenta corrigir com Hamming, depois verifica o EDC.

  \begin{lstlisting}[language=Python]
bloco = ADICIONAR_EDC[config["deteccao"]](bloco)
if config["correcao"] == "hamming":
    bloco = codificar_hamming(bloco)
  \end{lstlisting}

  \item[\rank{10}.] \textbf{Como vocês testaram o projeto e como provaram que o CRC-32 está correto?}

  O arquivo \texttt{testes.py} testa conversão texto/bits, enquadramentos, paridade, checksum, CRC, Hamming, modulações banda-base, modulações por portadora, ruído e simulação completa.

  O CRC-32 é validado com o vetor clássico:

  \begin{quote}
  \texttt{123456789 -> 0xCBF43926}
  \end{quote}

  \item[\rank{11}.] \textbf{Como funciona o enquadramento com flags, byte stuffing e bit stuffing?}

  No byte stuffing, o quadro começa e termina com \texttt{0x7E}. Se \texttt{0x7E} ou \texttt{0x7D} aparecem dentro do payload, o transmissor insere \texttt{0x7D} antes para marcar que aquilo é dado, não controle.

  No bit stuffing, a flag é \texttt{01111110}. Para essa flag não aparecer dentro do payload, o transmissor insere um zero depois de cinco bits 1 consecutivos.

  \item[\rank{12}.] \textbf{Qual é a diferença entre paridade, checksum e CRC-32?}

  Paridade é simples e fraca: detecta quantidade ímpar de erros. Checksum soma palavras de 16 bits em complemento de 1 e é intermediário. CRC-32 usa divisão polinomial em base 2 e é o detector mais forte do projeto.

  \begin{quote}
  Robustez: paridade < checksum < CRC-32. Custo no simulador: paridade 1 byte, checksum 2 bytes, CRC-32 4 bytes.
  \end{quote}

  \item[\rank{13}.] \textbf{Por que 16-QAM é mais sensível ao ruído do que QPSK?}

  Porque 16-QAM tem 16 pontos na constelação, enquanto QPSK tem 4. Como há mais pontos no mesmo espaço de sinal, eles ficam mais próximos. Com ruído, o ponto recebido pode ser deslocado e o receptor pode escolher o vizinho errado.

  \begin{quote}
  16-QAM ganha eficiência, mas perde margem contra ruído.
  \end{quote}

  \item[\rank{14}.] \textbf{O que significa constelação I/Q e para que serve a função de correlação?}

  Constelação I/Q é representar cada símbolo por duas coordenadas: I, em fase com o cosseno, e Q, em quadratura com o seno. Na transmissão, o código gera a onda usando essas coordenadas. Na recepção, a correlação estima qual par I/Q estava no sinal recebido.

  \begin{lstlisting}[language=Python]
soma_cos += amostra * math.cos(angulo)
soma_sen += amostra * math.sin(angulo)
i = 2 / N * soma_cos
q = -2 / N * soma_sen
  \end{lstlisting}

  \item[\rank{15}.] \textbf{Como ASK, FSK, QPSK e 16-QAM funcionam especificamente no código?}

  ASK muda amplitude: bit 1 transmite portadora, bit 0 transmite amplitude zero. FSK muda frequência: bit 0 usa uma frequência e bit 1 usa outra. QPSK agrupa 2 bits por símbolo e escolhe um ponto I/Q. 16-QAM agrupa 4 bits por símbolo: 2 bits escolhem I e 2 bits escolhem Q.

  \begin{quote}
  ASK e FSK carregam 1 bit por símbolo. QPSK carrega 2. 16-QAM carrega 4.
  \end{quote}

  \item[\rank{16}.] \textbf{O que acontece se o ruído aumentar muito?}

  O demodulador começa a errar bits. O Hamming pode corrigir erro simples por bloco, mas erros múltiplos podem ser apenas detectados ou até escapar, dependendo do EDC usado.

  \begin{quote}
  Quanto maior o ruído, maior a chance de o sinal recebido cair do lado errado da decisão.
  \end{quote}

  \item[\rank{17}.] \textbf{Por que o código completa bits com zero em QPSK e 16-QAM? Isso altera a mensagem?}

  QPSK precisa de grupos de 2 bits e 16-QAM precisa de grupos de 4 bits. Se o último grupo não fecha, o código completa com zeros só para formar o último símbolo. Isso não altera a mensagem útil, porque a camada de enlace e a aplicação recuperam apenas os bits pertencentes aos quadros válidos.

  \begin{lstlisting}[language=Python]
def pad(bits, tamanho):
    resto = len(bits) % tamanho
    if resto == 0:
        return bits
    faltam = tamanho - resto
    return bits + [0] * faltam
  \end{lstlisting}

  \item[\rank{18}.] \textbf{Como o receptor decide se recebeu 0 ou 1 nas modulações banda-base?}

  No NRZ-Polar, o receptor calcula a média das amostras do bit: média positiva vira 1, caso contrário vira 0. No Manchester, compara a média da primeira metade com a segunda: se começou mais alto e terminou mais baixo, é 1; se começou baixo e terminou alto, é 0. No Bipolar AMI, olha o módulo da média: se há pulso, é 1; se está perto de zero, é 0.

  \item[\rank{19}.] \textbf{Por que o projeto separa camada física, meio de comunicação e camada de enlace em arquivos diferentes?}

  Porque cada arquivo representa uma responsabilidade da pilha. A camada de enlace cria quadros, adiciona detecção e correção. A camada física transforma bits em sinal e sinal em bits. O meio de comunicação simula o canal, adicionando ruído. Essa separação deixa o código mais fiel ao modelo de redes e facilita explicar onde cada fenômeno acontece.

  \item[\rank{20}.] \textbf{Se o professor pedir para abrir o código, quais funções você deve saber explicar primeiro?}

  As funções mais importantes são:

  \begin{itemize}
    \item \texttt{meio\_comunicacao.transmitir}: mostra onde o ruído entra;
    \item \texttt{camada\_fisica.onda} e \texttt{correlacionar}: explicam portadora e I/Q;
    \item \texttt{modular\_manchester}: explica sincronismo;
    \item \texttt{adicionar\_paridade\_par}: explica paridade por quadro;
    \item \texttt{calcular\_crc32}: explica CRC manual;
    \item \texttt{codificar\_hamming} e \texttt{decodificar\_hamming}: explicam correção;
    \item \texttt{transmitir} e \texttt{receber} da camada de enlace: explicam a ordem do pipeline.
  \end{itemize}
\end{enumerate}

\end{document}
